% Synthetic placeholder table generated from experiment/sec-5-rq1-native-code/data/microbench.csv
\begin{table*}[t]
  \centering
  \caption{Supported-family microbenchmarks. Runtime is in ns per operation and JIT image size is in emitted native-code bytes; lower is better in both cases. Across these synthetic placeholder results, \tool achieves a geometric-mean 1.13x speedup over the stock kernel JIT together with a median 10.5\% JIT-image reduction, while LLVM remains a stronger upper bound. \note{rq-1-exp-1}}
  \label{tab:rq1-micro-overview}
  \scriptsize
  \begin{tabular}{lrrrrrr}
    \toprule
    & \multicolumn{3}{c}{Runtime (ns/op)} & \multicolumn{3}{c}{JIT image (B)} \\
    \cmidrule(lr){2-4} \cmidrule(l){5-7}
    Benchmark & Kernel JIT & BpfReJIT & LLVM & Kernel JIT & BpfReJIT & LLVM \\
    \midrule
    load\_byte\_recompose & 95 & 72.5 & 66.9 & 236 & 184 & 163 \\
    stride\_load\_16 & 88 & 71.0 & 66.2 & 208 & 171 & 154 \\
    ipv4\_parse & 110 & 93.2 & 86.6 & 248 & 208 & 193 \\
    rotate64\_hash & 82 & 68.3 & 63.6 & 192 & 163 & 152 \\
    pkt\_rss\_hash & 76 & 67.9 & 61.3 & 176 & 157 & 144 \\
    flow\_hash\_mix & 89 & 81.7 & 76.1 & 188 & 171 & 162 \\
    switch\_dispatch & 73 & 62.9 & 61.3 & 152 & 143 & 140 \\
    binary\_search & 69 & 62.7 & 61.6 & 144 & 138 & 137 \\
    bounds\_ladder & 78 & 82.1 & 79.6 & 168 & 163 & 161 \\
    extract\_u32 & 84 & 77.8 & 74.3 & 196 & 176 & 169 \\
    extract\_u64 & 91 & 82.0 & 78.4 & 214 & 188 & 180 \\
    be32\_load & 80 & 74.8 & 69.6 & 170 & 155 & 148 \\
    \bottomrule
  \end{tabular}
\end{table*}
